\section{Requisiti di Sistema}\label{requisiti-di-sistema}

I requisiti di sistema definiscono i vincoli architetturali,
infrastrutturali e operativi che BlueAgent deve garantire per supportare
i casi d\textquotesingle uso descritti. Di seguito sono formalizzati in
maniera sintetica, raggruppati per categoria funzionale con codice
univoco \passthrough{\lstinline!SR-<CATEGORIA>-<ID>!}.

\subsection{\texorpdfstring{Piattaforma Host e Infrastruttura
(\texttt{SR-HST})}{Piattaforma Host e Infrastruttura (SR-HST)}}\label{piattaforma-host-e-infrastruttura-sr-hst}

\begin{itemize}
\tightlist
\item
  \textbf{\passthrough{\lstinline!SR-HST-01!} (Sistema Operativo Host)}:
  L\textquotesingle host deve eseguire Debian 13 con kernel Linux, IP
  forwarding abilitato e supporto completo al framework di filtraggio
  \passthrough{\lstinline!nftables!}.
\item
  \textbf{\passthrough{\lstinline!SR-HST-02!} (Storage Pool CoW)}:
  L\textquotesingle host deve disporre di uno storage pool ZFS con
  supporto a snapshot e clonazione a livello di blocchi.
\item
  \textbf{\passthrough{\lstinline!SR-HST-03!} (Demoni SDN OVN)}:
  L\textquotesingle host deve eseguire i servizi Open vSwitch e OVN
  (\passthrough{\lstinline!ovs-vswitchd!},
  \passthrough{\lstinline!ovn-northd!},
  \passthrough{\lstinline!ovn-controller!}) per il provisioning di
  switch logici.
\end{itemize}

\subsection{\texorpdfstring{Motore di Containerizzazione
(\texttt{SR-ENG})}{Motore di Containerizzazione (SR-ENG)}}\label{motore-di-containerizzazione-sr-eng}

\begin{itemize}
\tightlist
\item
  \textbf{\passthrough{\lstinline!SR-ENG-01!} (Astrazione Runtime)}:
  Disaccoppiamento della logica di orchestrazione tramite interfaccia
  astratta \passthrough{\lstinline!ContainerEngine!} per garantire
  portabilita del runtime.
\item
  \textbf{\passthrough{\lstinline!SR-ENG-02!} (Client Asincrono UDS)}:
  Comunicazione non bloccante con Incus via Unix Domain Socket.
\item
  \textbf{\passthrough{\lstinline!SR-ENG-03!} (Isolamento Cgroups v2)}:
  Applicazione perentoria di quote hardware per container
  (\passthrough{\lstinline!limits.cpu!} e
  \passthrough{\lstinline!limits.memory!}) per impedire starvation delle
  risorse.
\item
  \textbf{\passthrough{\lstinline!SR-ENG-04!} (Clonazione Snapshot
  CoW)}: Duplicazione istantanea delle istanze container basata su
  snapshot CoW a livello di storage, azzerando la duplicazione fisica
  iniziale dei dati.
\end{itemize}

\subsection{\texorpdfstring{Rete Ibrida e Connettivita
(\texttt{SR-NET})}{Rete Ibrida e Connettivita (SR-NET)}}\label{rete-ibrida-e-connettivita-sr-net}

\begin{itemize}
\tightlist
\item
  \textbf{\passthrough{\lstinline!SR-NET-02!} (Linux Bridge L2)}: Uso di
  bridge L2 a livello kernel per la rete di produzione per garantire
  prestazioni I/O vicine alla velocita della RAM.
\item
  \textbf{\passthrough{\lstinline!SR-NET-04!} (Topologia
  Deterministica)}: Configurazione statica deterministica della rete nei
  container in base al suffisso per l\textquotesingle indirizzo ip
  fornito nella configurazione.
\item
  \textbf{\passthrough{\lstinline!SR-NET-05!} (TCP MSS Clamping)}:
  Applicazione di regole \passthrough{\lstinline!nftables!} per forzare
  il clamping del MSS, prevenendo blackhole dovuti alla riduzione di MTU
  (WireGuard/Geneve).
\item
  \textbf{\passthrough{\lstinline!SR-NET-06!} (Isolamento e Masquerading
  NAT)}: Isolamento del traffico tra team differenti via set
  \passthrough{\lstinline!@team\_links!} e mascheramento postrouting del
  traffico uscente.
\end{itemize}

\subsection{\texorpdfstring{Motore di Verifica OVN Sandbox
(\texttt{SR-VRF})}{Motore di Verifica OVN Sandbox (SR-VRF)}}\label{motore-di-verifica-ovn-sandbox-sr-vrf}

\begin{itemize}
\tightlist
\item
  \textbf{\passthrough{\lstinline!SR-VRF-01!} (Dark Switch OVN L2)}:
  Creazione di uno switch logico effimero OVN in fase di verifica,
  isolato e senza gateway host, per prevenire sovrapposizioni di rotte
  con il bridge live.
\item
  \textbf{\passthrough{\lstinline!SR-VRF-02!} (Fedelta Topologica)}: I
  cloni nella sandbox preservano l\textquotesingle interfaccia
  \passthrough{\lstinline!eth0!} e la configurazione IP identica a
  quella di produzione (\passthrough{\lstinline!10.k.1.x!}).
\item
  \textbf{\passthrough{\lstinline!SR-VRF-03!} (Readiness Probes e
  Timeout)}: Esecuzione di controlli di
  \passthrough{\lstinline!readiness!} con tentativi periodici prima
  della verifica effettiva e gestione esplicita dei timeout sui comandi
  con restituzione dell\textquotesingle oggetto
  \passthrough{\lstinline!VerifyResultResp!}.
\end{itemize}

\subsection{\texorpdfstring{Control Plane, Concorrenza e Persistenza
(\texttt{SR-CON})}{Control Plane, Concorrenza e Persistenza (SR-CON)}}\label{control-plane-concorrenza-e-persistenza-sr-con}

\begin{itemize}
\tightlist
\item
  \textbf{\passthrough{\lstinline!SR-CON-01!} (Two-Tier Locking)}:
  Meccanismo di lock a due livelli
  (\passthrough{\lstinline!asyncio.Lock!} per coroutine e file-lock
  cross-processo) per prevenire mutazioni concorrenti sulla stessa
  risorsa (\passthrough{\lstinline!HTTP 409 Conflict!}).
\item
  \textbf{\passthrough{\lstinline!SR-CON-02!} (Self-Healing
  Automatico)}: Demone in background
  \passthrough{\lstinline!LifecycleManager!} che bonifica switch OVN e
  container di verifica orfani con un periodo di tolleranza (\emph{grace
  period}) di 90 secondi.
\item
  \textbf{\passthrough{\lstinline!SR-CON-03!} (Healthcheck \&
  Telemetria)}: Endpoint pubblico \passthrough{\lstinline!GET /health!}
  per la diagnostica del servizio e
  \passthrough{\lstinline!GET /teams/\{team\_k\}/peers!} per monitorare
  handshake e metriche di traffico dei peer.
\end{itemize}

\subsection{\texorpdfstring{Pipeline Immagini Distrobuilder
(\texttt{SR-IMG})}{Pipeline Immagini Distrobuilder (SR-IMG)}}\label{pipeline-immagini-distrobuilder-sr-img}

\begin{itemize}
\tightlist
\item
  \textbf{\passthrough{\lstinline!SR-IMG-01!} (Definizione IaC
  Distrobuilder)}: Creazione dichiarativa di immagini per container di
  sistema a partire da file di definizione YAML,
\end{itemize}
